\documentclass[11pt]{article}
% Спільна преамбула для звітів (XeLaTeX)
\usepackage{fontspec}
\usepackage{polyglossia}
\setmainlanguage{ukrainian}
\setotherlanguage{english}
\setmainfont{DejaVuSerif}[Extension=.ttf,UprightFont=*,BoldFont=*-Bold,ItalicFont=*-Italic,BoldItalicFont=*-BoldItalic]
\setsansfont{DejaVuSans}[Extension=.ttf,UprightFont=*,BoldFont=*-Bold,ItalicFont=*-Oblique,BoldItalicFont=*-BoldOblique]
\setmonofont{DejaVuSansMono}[Extension=.ttf,UprightFont=*,BoldFont=*-Bold,ItalicFont=*-Oblique,BoldItalicFont=*-BoldOblique,Scale=0.9]
\usepackage[a4paper,margin=2.2cm]{geometry}
\usepackage{booktabs}
\usepackage{tabularx}
\usepackage{array}
\usepackage{enumitem}
\usepackage{xcolor}
\usepackage{tcolorbox}
\usepackage{graphicx}
\usepackage{setspace}
\usepackage[hidelinks]{hyperref}
\usepackage{titlesec}
\usepackage{fancyhdr}
\usepackage{multirow}
\usepackage{colortbl}
\setlength{\parskip}{4pt}
\setlength{\parindent}{0pt}
\onehalfspacing

\definecolor{accent}{HTML}{1F4E79}
\definecolor{sgreen}{HTML}{E2F0D9}
\definecolor{wred}{HTML}{FBE5D6}
\definecolor{oblue}{HTML}{DEEBF7}
\definecolor{tyellow}{HTML}{FFF2CC}
\definecolor{lightgray}{HTML}{F2F2F2}

\titleformat{\section}{\large\bfseries\color{accent}}{\thesection}{0.6em}{}
\titleformat{\subsection}{\normalsize\bfseries}{\thesubsection}{0.6em}{}

\pagestyle{fancy}
\fancyhf{}
\lhead{\small ІТ, Модуль 2 — командний проєкт}
\rhead{\small Петров Б., Мекшаков А.}
\cfoot{\thepage}
\renewcommand{\headrulewidth}{0.3pt}

\newcolumntype{L}[1]{>{\raggedright\arraybackslash}p{#1}}
\newcolumntype{Y}{>{\raggedright\arraybackslash}X}

% Титульна сторінка
\newcommand{\maketitlepage}[3]{%
\begin{titlepage}
\centering
{\large Київський національний університет імені Тараса Шевченка\\[2pt]
Факультет комп'ютерних наук та кібернетики\par}
\vspace{3cm}
{\large Дисципліна: \textbf{Інформаційні технології}\par}
\vspace{0.5cm}
{\large Модуль 2 — робота в команді\par}
\vspace{2cm}
{\LARGE\bfseries #1\par}
\vspace{0.6cm}
{\large #2\par}
\vfill
\begin{flushright}
\begin{tabular}{l}
\textbf{Виконали:} студенти 4 курсу, група МІ-41\\[4pt]
Петров Богдан\\
Мекшаков Артемій\\[10pt]
\textbf{Спеціальність:} Комп'ютерні науки\\
\end{tabular}
\end{flushright}
\vspace{1.5cm}
{Київ — #3\par}
\end{titlepage}
}

\geometry{margin=2cm}
\setlength{\parskip}{3pt}
\setstretch{1.12}
\newcommand{\say}[1]{\begin{quote}\itshape «#1»\end{quote}}
\newcommand{\who}[1]{\textcolor{accent}{\textbf{[#1]}}}
\newcommand{\todo}[1]{\textcolor{red}{\textbf{[TODO: #1]}}}
\begin{document}
\setlength{\emergencystretch}{2em}

\begin{center}
{\LARGE\bfseries Сценарій захисту етапу 4}\\[3pt]
{\small $\approx$ 8 хв. Б = Богдан, А = Артемій. Відкрити заздалегідь: PDF презентації; браузер --- звичайне вікно (\texttt{starosta@studyflow.dev}) і інкогніто (\texttt{artem@studyflow.dev}), пароль \texttt{Demo1234!}; Swagger \texttt{localhost:3001/api/docs}; Mailpit \texttt{localhost:8025}; GitHub-репо, вкладка Actions (зелений CI); Notion --- вид «Спринти». Перед парою: \texttt{docker compose up -d}, \texttt{pnpm db:reset}, \texttt{pnpm dev}, очистити Mailpit.}
\end{center}

\who{Б} \textbf{Слайд 2 --- StudyFlow.}
\say{Нагадаю ідею: староста створює дедлайн один раз --- він з'являється в календарі кожного, кожен відмічає виконання сам, а група бачить прогрес. За етап зробили робочий прототип на Next.js, NestJS і PostgreSQL з тестами, Swagger і UML.}

\who{А} \textbf{Слайд 3 --- План спринта.}
\say{Обсяг етапу --- спринти 1--3: 31 елемент, 68 story points. Усі вони в Done --- Must, Should і навіть Could, перетягування в канбані. Статуси в Notion відповідають репозиторію.}

\who{Б} \textbf{Слайд 4 --- Архітектура.}
\say{Модульний моноліт: один NestJS-застосунок, модулі за доменом --- auth, групи, предмети, дедлайни, задачі, нагадування --- і шари контролер, сервіс, Prisma. Фронтенд звертається до /api на тому ж домені, Next.js проксіює запити, тож refresh-cookie працює без CORS. Типи для фронтенду генеруються зі Swagger.}

\who{Б} \textbf{Слайд 5 --- Модель даних.}
\say{Вісім таблиць. Груповий і особистий дедлайн --- одна таблиця: у особистого group\_id порожній. Відмітки виконання --- окрема таблиця з ключем «дедлайн --- користувач», тому прогрес групи --- це просто кількість відміток поділена на кількість учасників.}

\who{Б} \textbf{Слайд 6 --- UML.}
\say{Зліва ролі: староста має все, що студент, плюс керування групою. Справа ключовий сценарій: guard перевіряє роль старости в БД, сервіс перевіряє, що предмет з цієї групи, учасник відмічає виконання ідемпотентно, і староста бачить 50 відсотків.}

\who{А} \textbf{Слайд 7 --- Демо.} Перемкнутись на браузер.
\begin{enumerate}[leftmargin=2.2em,itemsep=1pt,topsep=3pt]
  \item Вікно старости: «Групи» $\to$ «МІ-41» $\to$ «Додати дедлайн»; назва «ЛР 6: демо на захисті», предмет «Інформаційні технології», термін уже стоїть --- завтра, 23:59 $\to$ «Додати дедлайн». Прогрес --- 0\,\%.
  \item Інкогніто (\texttt{artem}): «Календар» $\to$ «Тиждень» --- дедлайн уже на завтра $\to$ клік по ньому $\to$ «Відмітити виконаним».
  \item Вікно старости: оновити сторінку групи --- заповнилась 1 клітинка з 4, прогрес 25\,\%.
  \item «Дошка задач»: перетягнути «Календар тижня» з «В роботі» в «Готово».
  \item Swagger, тег \texttt{dev}: \texttt{POST /api/dev/reminders/run} $\to$ Try it out $\to$ Execute; \texttt{sent} у відповіді --- скільки листів пішло $\to$ Mailpit, листи-нагадування.
  \item Swagger: тег \texttt{deadlines}, схема \texttt{CreateDeadlineDto}.
\end{enumerate}
\say{Староста створює дедлайн один раз --- і учасник одразу бачить його у своєму календарі. Він відмічає виконання, а староста бачить прогрес групи. Нагадування приходять лише тим, хто ще не відмітив.}

\who{А} \textbf{Слайди 8--9 --- Скріншоти.} Лише якщо живе демо не піде.

\who{Б} \textbf{Слайд 10 --- Модульні тести.}
\say{Сорок модульних тестів, усі проходять. Prisma замокано, тому кожен тест перевіряє одне бізнес-правило: дубль email, ротацію refresh-токена, останнього старосту, прогрес, повторні нагадування. ID тестів збігаються з тест-планом, журнал скрипт будує сам.}

\who{А} \textbf{Слайд 11 --- Інтеграційні та E2E.}
\say{Тринадцять інтеграційних тестів ганяють увесь застосунок проти справжньої тестової бази, пошта --- заглушка. Три E2E у Playwright, головний --- сценарій із Definition of Done у двох браузерах: реєстрація, група, дедлайн, вступ за кодом, відмітка, 50 відсотків.}

\who{Б} \textbf{Слайд 12 --- API.}
\say{34 операції у восьми тегах, кожна з описом, кодами помилок і прикладами DTO. Специфікація лежить у репозиторії, з неї генеруються типи фронтенду --- якщо API зміниться, фронтенд не збереться.}

\who{А} \textbf{Слайд 13 --- Рефлексія.}
\say{Успіх --- увесь обсяг і тестів більше за план. Неуспіх --- реалізацію почали лише 10 листопада: спринти 1--2 пішли на документацію етапів 2--3, а дедлайн етапу~4 перенесли на 15-те. Працювали парно з одного ноутбука, тому коміти з одного акаунта і без рев'ю, а CI у GitHub запустився лише на готовому коді. У спринті~4: кожна задача в гілці з власного акаунта, pull request з рев'ю і CI, деплой --- першим.}

\who{Б} \textbf{Слайд 14 --- Наступні кроки.}
\say{До першого грудня: деплой на Vercel і Railway, налаштування нагадувань і особиста статистика, пілот у нашій групі і фінальний звіт. Дякуємо.}

\section*{Якщо спитають}
\begin{itemize}[leftmargin=1.2em,itemsep=2pt]
  \item \textbf{Чому моноліт, а не мікросервіси?} --- Двоє розробників і один домен; мікросервіси додали б лише інфраструктуру. Модулі вже розділені --- винести можна пізніше.
  \item \textbf{Чому Prisma, а не TypeORM?} --- Схема --- єдине джерело правди: з неї міграції, типи запитів і ER-діаграма.
  \item \textbf{Де токени?} --- Access на 15 хв у пам'яті (не в localStorage --- XSS), refresh на 7 днів у httpOnly cookie з ротацією, у БД лише SHA-256-хеш; logout його стирає.
  \item \textbf{Ролі, а якщо обійти UI?} --- GroupRoleGuard на сервері бере роль із БД на кожен запит; запит учасника на дію старости --- 403, це перевіряє IT-06.
  \item \textbf{Що мокається?} --- У модульних --- Prisma, сервіс без DI. В інтеграційних --- справжня PostgreSQL \texttt{studyflow\_test}, заглушка лише пошта.
  \item \textbf{Як без дублікатів листів?} --- Таблиця reminder\_logs з ключем «дедлайн --- користувач»: пару з записом пропускаємо. IT-12: 2 листи, повторний запуск --- 0.
  \item \textbf{Coverage і чому не 100\,\%?} --- 54\,\% рядків модульними, 73{,}5\,\% інтеграційними. Модульні --- про бізнес-правила, контролери перевіряють інтеграційні; 100\,\% не мета.
  \item \textbf{Що не встигли?} --- З обсягу спринтів 1--3 --- нічого. Не встигли процес: PR із рев'ю, а CI у GitHub уперше запустився лише 11.11, на готовому коді (зелений: 40/40 і 13/13). Деплой --- етап~5, за планом.
\end{itemize}

\textbf{Цифри:} 31 ID / 68 SP, усе Done · тести 40 / 13 / 3 · coverage 54\,\% / 73{,}5\,\% · 34 операції API, 8 тегів · 8 таблиць, 7 UML · CI у GitHub зелений.

\end{document}
